\documentclass{article}
\usepackage[T1]{fontenc}
\usepackage{lmodern}
\usepackage{graphicx}
\usepackage{amsmath,amssymb}
\usepackage[a4paper,margin=2cm]{geometry}
\usepackage[absolute,overlay]{textpos}
\setlength{\TPHorizModule}{1mm}
\setlength{\TPVertModule}{1mm}
\usepackage[indonesian]{babel}
\usepackage{hyperref}
\setlength{\emergencystretch}{2em}
% unit: Halaman 73

\begin{document}

% === Halaman 73 (python/native) ===

73

Contoh  plainteks:  temui ibu nanti malam 

   → Tidak ada huruf j, maka langsung tulis  pesan dalam pasangan huruf

        (setelah semua spasi dibuang):


te mu ii bu na nt im al am 


 → Ada bigram dengan huruf yang berulang (ii), sisipkan huruf x di tengahnya:  


te mu ix ib un an ti ma la m  


  → Tambahkan huruf x jika bigram terakhir hanya satu huruf:


te mu ix ib un an ti ma la mx

\end{document}
